\section{INTRODUCTION}\label{sec:intro}

\IEEEPARstart{O}{n-orbit} servicing and the removal of failed spacecraft require a servicing spacecraft to approach a target autonomously and to establish the relative position, attitude, and velocity needed for capture \cite{tian2024taes}. In this pre-capture phase, the chaser must respect a safe distance, a narrow approach corridor, the field of view of its camera, relative-speed limits, and its actuator capability, so that a capture window opens \cite{lian2026asr}. A noncooperative target, such as a failed spacecraft \cite{xu2026ast}, offers no assistance: its docking interface moves with the tumbling body, and the chaser must coordinate its relative position and attitude with a moving capture point rather than a fixed one \cite{dong2025aa,wang2026ast}. Fast tumbling turns this into a timing problem. When the target tumbles at $2$--$3^{\circ}$/s, the port faces the chaser only once per spin period of about two and a half minutes; with a deadline and a propellant budget, this period is comparable to the deadline. The feasible capture window then varies with time and is scarce: a missed pass costs a full spin period, and whether the next pass is usable depends on the approach speed and timing chosen tens of seconds before it.

Constrained approaches to tumbling targets are commonly solved by model predictive control (MPC) or trajectory optimization. A two-stage MPC brings the chaser into a time-varying line-of-sight region and then tracks the periodic motion of the port \cite{dong2022aa}; tube MPC with online inertia estimation accounts for the limited predictability of the tumbling \cite{dong2024asr}; a dual-layer MPC filters nominal commands with control barrier functions \cite{xu2026ast}; second-order cone MPC encodes corridor and pointing requirements as convex constraints \cite{lian2026asr}; and manifold-based planning computes soft-docking trajectories \cite{wang2026ast}. These methods enforce the constraints explicitly over their prediction horizons and reach the port by converging to its periodic motion or by tracking it within one alignment. They consider moderate tumbling rates, such as about $1^{\circ}$/s \cite{tian2024taes}, or tasks without a deadline and budget \cite{dong2024asr}, in which a missed alignment can be awaited, and their long-horizon performance relies on the reference and terminal design \cite{xu2026ast}. Which pass to use, and how to arrive at it in time, lies beyond the horizon of such an optimization.

Reinforcement learning (RL), in contrast, learns state-dependent policies from interaction with the objective of the whole task. Deep RL algorithms such as soft actor--critic (SAC) \cite{haarnoja2018sac} and TD3 \cite{fujimoto2018td3} have controlled six-degree-of-freedom rendezvous with a rotating target end to end \cite{yang2025taes}, and hierarchical and safe RL have been applied to small-satellite guidance and to attitude tracking of noncooperative targets \cite{tammam2025cep,lu2026ast}; under complex constraints, however, exploration efficiency and execution reliability remain limited. MPC and RL thus have nearly complementary properties \cite{reiter2026arc,reiter2025ac4mpc}, and their combinations can be grouped by what the learned component provides to the optimizer \cite{reiter2026arc}. Learning can supply value and cost information: the terminal value, the cost, or the initial guess of an MPC \cite{reiter2025ac4mpc,ghezzi2026ecc,gong2026ast}, value targets from model predictions \cite{wu2024mpcvalue}, or MPC-shaped rewards for docking \cite{ramezani2026aa}. Learning can also supply task references and timing to an MPC that performs the constrained execution, such as maneuver times \cite{ren2026taes} and re-optimization instants \cite{zheng2026aa}, or coordinate the authority of learned and model-based controllers by flight phase \cite{adjei2026ast}. This article belongs to the second group: the learned layer sets stage targets, and the executor is part of its environment \cite{reiter2026arc}. In such schemes a learned decision is executed as issued, and its consequence reaches learning only through the return; where the executor may change a decision, it does so for safety or feasibility, as in shields \cite{alshiekh2018shielding}, feasibility and reference governors \cite{leung2024scl,tartaglione2024aut,duan2025autcon}, and safe-evolving or residual RL \cite{lu2025eaai,zeng2026neucom}.

Fast-tumbling pre-capture under a deadline leaves two problems open in these architectures. The first is the use of the dynamic capture window. Whether a pass is usable depends on how closely the line of sight approaches the cone swept by the port, and when it occurs depends on the phase of the port; a locally feasible approach does not imply that a pass is reached before the deadline. The decisions of the learned layer must therefore act on these quantities. The second is the coordination between long-horizon decisions and short-horizon constrained execution. Whether a decision is converted into an entry is resolved by the planning, gating, and tracking of the executor within seconds, whereas the decision layer acts once per stage and the outcome, capture or timeout, appears much later. When passes are scarce, neighboring decisions lead to capture and to timeout, so that the value of a decision changes abruptly across the action space at boundaries set by the executor. Such a value is hard to learn from the sparse terminal return, and a modification aimed at safety or feasibility leaves a safe decision unchanged even when it misses the last usable opportunity.

This article addresses the two problems with a hierarchical framework in which a learned policy issues two-dimensional stage intents to an execution layer of planners and MPC. Rotation-synchronized window guidance (RWG) turns each intent into stage goals that regulate the line of sight onto the sweep cone of the port and set the closing pace and the timing of the encounter, and it plans and re-times entries toward the predicted port motion. Execution-feedback closed-loop coordination (ECC) lets the execution layer predict, at every decision, the entry opportunities that it can realize, selects the executed object from the proposal and these opportunities by a fixed rule, and trains the policy on the resulting closed loop. The main contributions are as follows.
\begin{enumerate}
\item An execution-opportunity coordination of learning and constrained execution. At every decision, the execution layer predicts with its own models whether the proposed intent and a grid of alternatives lead to an entry, and whether an entry can still be planned directly at the last usable pass; one rule keeps the proposal, corrects it at a critical window, or launches the direct entry. SAC is trained on the coordinated closed loop with the proposals stored, so that the critics evaluate what happens when the policy proposes an intent, with the losses of standard SAC. In contrast to shields and governors \cite{alshiekh2018shielding,leung2024scl,lu2025eaai}, the coordination acts on task opportunities rather than on safety or feasibility, and in contrast to model-based value estimation \cite{wu2024mpcvalue}, the execution layer generates no value targets.
\item A rotation-synchronized window guidance derived from the sweep-cone geometry of the port. It characterizes the passing window by the latitude offset and the phase of the line of sight, regulates the latitude by a fixed law, and leaves the closing pace and the along-sweep timing to the policy; entries are gated on the predicted arrival and re-timed within the same pass. Unlike multistage MPC that converges to the periodic motion of the port \cite{dong2022aa}, RWG schedules the encounter with a pass that may not recur before the deadline.
\item A validation in six-degree-of-freedom simulations of time-limited pre-capture with gravity gradient and $J_{2}$ perturbation, with safety judged on the true trajectory and the dynamic model of \cite{yang2025taes}. On paired test tasks, ECC captures $85.4\%$ of the tasks, $11.6$ percentage points more than learning-guided MPC on the same execution chain at comparable time and propellant cost, and it retains this capture rate under a slower spin and a stronger nutation outside the training distribution.
\end{enumerate}

The remainder of this article is organized as follows. Section~\ref{sec:problem} formulates the pre-capture problem. Section~\ref{sec:framework} presents the hierarchical framework and RWG. Section~\ref{sec:ecc} develops ECC. Section~\ref{sec:experiments} reports the numerical study, and Section~\ref{sec:conclusion} concludes the article.

\textit{Notations:} Bold lowercase and uppercase symbols denote vectors and matrices, $\norm{\cdot}$ is the Euclidean norm, $\norm{\vect{y}}_{\vect{W}}^{2}=\vect{y}^{\T}\vect{W}\vect{y}$, $\bI_{3}$ is the identity matrix, $\vect{1}$ is the vector of ones, $\vect{y}_{1}^{\times}$ is the skew-symmetric matrix with $\vect{y}_{1}^{\times}\vect{y}_{2}=\vect{y}_{1}\times\vect{y}_{2}$, $(\cdot)^{\wedge}$ maps a twist to $\mathfrak{se}(3)$, $\log$ is the matrix logarithm on $\SO$ and $(\cdot)^{\vee}$ the inverse of $(\cdot)^{\times}$, and $\angle(\vect{y}_{1},\vect{y}_{2})$ is the angle between two vectors. $\norm{\cdot}_{\infty}$ is the maximum norm, and the saturation $\sat(\vect{y};\bar y)=\vect{y}\min\{1,\bar y/\norm{\vect{y}}\}$ for $\vect{y}\neq\vect{0}$, with $\sat(\vect{0};\bar y)=\vect{0}$, applies to scalars and vectors. A hat denotes a predicted quantity, a bar an upper bound, and a tilde a transformed quantity; $\Delta$ denotes increments only.
